\section{Contexto e Problema}

\subsection{Contexto Atual}

A arquitetura atual da empresa está em processo de evolução de um modelo predominantemente monolítico para uma arquitetura híbrida, na qual o monólito passa a assumir um papel mais estratégico como centralizador de informações e orquestrador de fluxos, enquanto componentes especializados são introduzidos para desacoplar responsabilidades específicas.

Como parte dessa evolução, foi desenvolvido o sistema denominado Integrador, cuja função é atuar como gateway de integrações com fornecedores externos, especialmente locadoras de veículos. Esse novo componente tem como objetivo simplificar o desenvolvimento, manutenção e evolução das integrações, padronizando a comunicação com fornecedores, abstraindo complexidades específicas de cada integração e promovendo maior escalabilidade e governança técnica.

Historicamente, as integrações com locadoras eram implementadas diretamente no monólito, o que resultava em um alto acoplamento, dificuldade de manutenção, baixa reutilização de componentes e maior complexidade para evolução contínua. O Integrador surge, portanto, como uma solução estruturante para resolver essas limitações ao longo do tempo.

Atualmente, o Integrador já se encontra em produção e suporta diversas integrações ativas, operando de forma estável e agregando valor ao ecossistema. No entanto, ainda existem integrações legadas que permanecem no monólito e precisam ser migradas para o novo modelo, consolidando a arquitetura proposta.

\subsection{Problemas Observados}

Apesar dos avanços proporcionados pelo Integrador, o processo de migração das integrações existentes apresenta desafios relevantes, especialmente em dois aspectos principais:

\begin{itemize}
   \item Complexidade das regras de integração:
   \begin{itemize}
     \item Parte significativa das integrações existentes possui lógica de negócio complexa e está encapsulada em drivers específicos dentro do monólito. Essas regras envolvem particularidades de fornecedores, tratamentos específicos de dados, fluxos condicionais e exceções que dificultam a replicação fiel do comportamento no novo gateway. Atualmente, o processo de migração é realizado por meio de scripts pré-definidos, que são constantemente ajustados conforme novas necessidades surgem, o que evidencia a ausência de um processo estruturado e previsível.
   \end{itemize}
   \item Ausência de um mecanismo robusto de validação de integração:
   \begin{itemize}
     \item O principal desafio identificado está na validação da equivalência entre o comportamento do sistema legado e o novo gateway durante a migração. Atualmente, essa validação depende majoritariamente de testes manuais, análises pontuais e monitoramento reativo em produção. Essa abordagem limita a capacidade de identificar divergências de forma sistemática, reduz a confiança no processo de migração e aumenta o risco de regressões funcionais.
   \end{itemize}
 \end{itemize}





\subsection{Impacto no Negócio}

A ausência de um processo estruturado de validação durante a migração gera impactos diretos e indiretos no negócio, dentre os quais se destacam:

\begin{itemize}
   \item Risco operacional elevado, com possibilidade de inconsistências em fluxos críticos como busca, precificação, disponibilidade e confirmação de reservas;
   \item Impacto na experiência do usuário, decorrente de falhas silenciosas ou divergências entre o comportamento esperado e o efetivamente entregue;
   \item Potencial perda de receita, associada a falhas em reservas, indisponibilidade de ofertas ou erros de integração com fornecedores;
   \item Baixa previsibilidade na migração, dificultando o planejamento e a priorização de integrações a serem migradas;
   \item Aumento do custo operacional, devido à necessidade de retrabalho, correções emergenciais e maior dependência de validações manuais.
\end{itemize}

Diante desse cenário, torna-se evidente a necessidade de um mecanismo estruturado que permita validar, de forma automatizada e confiável, a equivalência entre os sistemas durante o processo de migração, reduzindo riscos e aumentando a eficiência da transição arquitetural.